% =====================================================================
% Tema 3 — Polarización y cristales ópticos (Saleh cap. 6). Problemas 18 a 27.
% Convención de Saleh: e^{j\omega t}, \mathcal E = Re{E e^{j\omega t}}, onda viajera e^{-jkz}.
% =====================================================================

\textbf{Convención de todo el tema.} Se usa la de Saleh y Teich: dependencia temporal
$e^{j\omega t}$, campo real $\Ec=\mathrm{Re}\{\mathbf A\,e^{j\omega t}\}$, onda que avanza en $+z$
con factor espacial $e^{-jkz}$, vector de Jones $\J=(A_x,A_y)^T$ y fase relativa
$\varphi=\varphi_y-\varphi_x$ (Saleh, §6.1A, ecs.~6.1-1 a 6.1-5, pp.~199--200). Una onda es
\emph{circular derecha} si $\varphi=+\pi/2$ y gira en sentido horario para un observador que la
ve venir (Saleh, p.~201). Las matrices que más se usan son:
\[
\R(\theta)=\begin{bmatrix}\cos\theta&\sin\theta\\-\sin\theta&\cos\theta\end{bmatrix}\ (6.1\text{-}22),\qquad
\T_{\rm rot}(\theta)=\begin{bmatrix}\cos\theta&-\sin\theta\\ \sin\theta&\cos\theta\end{bmatrix}\ (6.1\text{-}20),\qquad
\T_{\rm ret}=\begin{bmatrix}1&0\\0&e^{-j\Gamma}\end{bmatrix}\ (6.1\text{-}19),
\]
y un elemento de matriz $\T'$ en ejes girados $\theta$ vale $\T=\R(-\theta)\T'\R(\theta)$ en los ejes
$x$--$y$ (ec.~6.1-24, p.~208). Para un retardador de retardo $\Gamma$ con eje rápido a $\theta$, con
$c=\cos\theta$, $s=\sin\theta$ y $e\equiv e^{-j\Gamma}$, el producto da
\begin{equation}
\T_{\rm ret}(\Gamma,\theta)=\R(-\theta)\begin{bmatrix}1&0\\0&e\end{bmatrix}\R(\theta)
=\begin{bmatrix}c^2+s^2e & cs(1-e)\\ cs(1-e) & s^2+c^2e\end{bmatrix}.
\label{eq:t3_retgen}
\end{equation}
(Se comprueba multiplicando: $\T'\R(\theta)=\left[\begin{smallmatrix}c&s\\-se&ce\end{smallmatrix}\right]$ y luego
$\R(-\theta)$ por la izquierda.)

%---------------------------------------------------------------------
\sol{18}{Elipse, Stokes y Poincaré}{Saleh, §6.1A, ecs.~6.1-6 a 6.1-9, pp.~200--203; Prob.~6.1-5, p.~240.}
\label{sol:18}

El vector es $\J=\frac{1}{\sqrt5}\,(1,\;2e^{j\pi/3})^T$, normalizado: $|A_x|^2+|A_y|^2=\frac{1}{5}+\frac45=1$.

\inc{a} Con $A_x=\mathsf a_xe^{j\varphi_x}$ y $A_y=\mathsf a_ye^{j\varphi_y}$:
$\mathsf a_x=1/\sqrt5$, $\varphi_x=0$, $\mathsf a_y=2/\sqrt5$, $\varphi_y=\pi/3$. Por tanto
\[
\mathsf r=\frac{\mathsf a_y}{\mathsf a_x}=2,\qquad \varphi=\varphi_y-\varphi_x=60^\circ .
\]
La orientación sale de (Saleh, ec.~6.1-6, p.~200):
\[
\tan2\psi=\frac{2\mathsf r}{1-\mathsf r^2}\cos\varphi=\frac{2\cdot2}{1-4}\cdot\frac12=-\frac23 .
\]
La tangente no fija el cuadrante: $2\psi=-33.69^\circ$ o $2\psi=146.31^\circ$. Para elegir se usan
los signos de $\cos2\psi$ y $\sin2\psi$, que son los de $\mathsf S_1$ y $\mathsf S_2$ (ecs.~6.1-8a,b, abajo):
$\mathsf S_1<0$ y $\mathsf S_2>0$, así que $2\psi$ está en el segundo cuadrante, $2\psi=146.31^\circ$.
Comprobación física: como $\mathsf a_y=2\mathsf a_x$, la elipse es alargada en $y$, y su eje mayor
debe estar más cerca del eje $y$ que del $x$; en efecto, $\psi=73.15^\circ$.
La elipticidad sale de (Saleh, ec.~6.1-7, p.~200):
\[
\sin2\chi=\frac{2\mathsf r}{1+\mathsf r^2}\sin\varphi=\frac{4}{5}\cdot\frac{\sqrt3}{2}=0.6928
\;\Rightarrow\;2\chi=43.85^\circ,\quad \chi=21.93^\circ .
\]
Los semiejes cumplen $b/a=\tan\chi=0.403$ y $a^2+b^2=1$, de donde $a=0.928$ y $b=0.374$.
\resultado{$\mathsf r=2$, $\varphi=60^\circ$, $\psi=73.2^\circ$, $\chi=21.9^\circ$.}

\inc{b} $0<\varphi<\pi$ (es decir, $\sin\varphi>0$ y $\chi>0$): la polarización es \textbf{derecha}.
Comprobación directa con el campo real en $z=0$, $\Ec_x=\frac{1}{\sqrt5}\cos\omega t$,
$\Ec_y=\frac{2}{\sqrt5}\cos(\omega t+60^\circ)$: en $\omega t=0$ el campo es
$(0.447,\,0.447)$; en $\omega t=30^\circ$ es $(0.387,\,0)$. El extremo baja del primer cuadrante al eje
$x$: gira en sentido horario visto con $x$ a la derecha, $y$ arriba y la onda viniendo hacia el lector.
Según Saleh (p.~201), eso es derecha.
\resultado{Elíptica derecha.}

\inc{c} Con las ecs.~6.1-9 (Saleh, p.~202), y $A_x^*A_y=\frac{2}{5}e^{j\pi/3}$:
\begin{align*}
\mathsf S_0&=|A_x|^2+|A_y|^2=1, &
\mathsf S_1&=|A_x|^2-|A_y|^2=\tfrac15-\tfrac45=-0.6,\\
\mathsf S_2&=2\,\mathrm{Re}\{A_x^*A_y\}=\tfrac45\cos60^\circ=0.4, &
\mathsf S_3&=2\,\mathrm{Im}\{A_x^*A_y\}=\tfrac45\sin60^\circ=0.6928 .
\end{align*}
Verificación: $\mathsf S_1^2+\mathsf S_2^2+\mathsf S_3^2=0.36+0.16+0.48=1.00=\mathsf S_0^2$. Es lo esperado para
luz totalmente polarizada (Saleh, p.~203).
Las ecs.~6.1-8 dan lo mismo: $\cos2\chi\cos2\psi=0.7211\cdot(-0.8321)=-0.600$ y
$\cos2\chi\sin2\psi=0.7211\cdot0.5547=0.400$.
\resultado{$\mathbf S=(1,\,-0.6,\,0.4,\,0.693)$; $\ \mathsf S_1^2+\mathsf S_2^2+\mathsf S_3^2=\mathsf S_0^2$.}

\inc{d} En la esfera de Poincaré el punto tiene coordenadas esféricas $(1,\ 90^\circ-2\chi,\ 2\psi)$
(Saleh, p.~201), es decir, ángulo polar $46.15^\circ$ y azimut $146.31^\circ$. Sus coordenadas
cartesianas son $(\mathsf S_1,\mathsf S_2,\mathsf S_3)/\mathsf S_0=(-0.6,\,0.4,\,0.693)$. Está en el
hemisferio norte (derecha), a $43.85^\circ$ de latitud, con azimut entre LP a $45^\circ$
($2\psi=90^\circ$) y LP$_y$ ($2\psi=180^\circ$), más cerca de LP$_y$. Ver la Fig.~\ref{fig:t3_p18}.

\inc{e} \emph{Estado ortogonal.} Se busca $\J_\perp$ con $(\J,\J_\perp)=A_xA_{\perp x}^*+A_yA_{\perp y}^*=0$
(Saleh, ec.~6.1-11, p.~203). Para cualquier $\J=(A_x,A_y)$ normalizado, el vector
$\J_\perp=(-A_y^*,\,A_x^*)$ cumple
$A_x(-A_y^*)^*+A_y(A_x^*)^*=-A_xA_y+A_yA_x=0$. Aquí
$\J_\perp=\frac{1}{\sqrt5}(-2e^{-j\pi/3},\,1)$; multiplicando por la fase global $-e^{j\pi/3}$,
que no cambia el estado,
\[
\J_\perp=\frac{1}{\sqrt5}\begin{bmatrix}2\\ -e^{j\pi/3}\end{bmatrix}=\frac{1}{\sqrt5}\begin{bmatrix}2\\ e^{-j2\pi/3}\end{bmatrix}.
\]
Entonces $\mathsf r_\perp=1/2$, $\varphi_\perp=-120^\circ$ y
$\tan2\psi_\perp=\frac{2(1/2)}{1-1/4}\cos(-120^\circ)=-\frac23$, con $\mathsf S_{1}=+0.6$ y
$\mathsf S_{2}=-0.4$. Por tanto $2\psi_\perp=-33.69^\circ$ y $\psi_\perp=-16.85^\circ\equiv163.15^\circ$.
La elipticidad es $\sin2\chi_\perp=\frac{2(1/2)}{5/4}\sin(-120^\circ)=-0.6928$, es decir,
$\chi_\perp=-21.93^\circ$. Se cumple:
\begin{itemize}
\item $\psi-\psi_\perp=73.15^\circ-(-16.85^\circ)=90^\circ$: los ejes mayores son perpendiculares;
\item $\chi_\perp=-\chi$: misma forma de elipse y sentido de giro opuesto (izquierda).
\end{itemize}
\emph{Demostración general (Prob.~6.1-5).} Con $\J_\perp=(-A_y^*,A_x^*)$:
$\mathsf S_{1\perp}=|A_y|^2-|A_x|^2=-\mathsf S_1$ y
$A_{\perp x}^*A_{\perp y}=(-A_y)(A_x^*)=-A_x^*A_y$, de modo que $\mathsf S_{2\perp}=-\mathsf S_2$ y
$\mathsf S_{3\perp}=-\mathsf S_3$. Dos estados ortogonales son \emph{antípodas} en la esfera. Las
ecs.~6.1-8 dicen entonces que $\sin2\chi_\perp=-\sin2\chi$, de donde $\chi_\perp=-\chi$ (giro opuesto), y que
$(\cos2\psi_\perp,\sin2\psi_\perp)=-(\cos2\psi,\sin2\psi)$, de donde $2\psi_\perp=2\psi\pm180^\circ$ y
$\psi_\perp=\psi\pm90^\circ$ (ejes mayores perpendiculares). $\blacksquare$
\resultado{$\J_\perp=\frac{1}{\sqrt5}(2,\,-e^{j\pi/3})^T$: $\psi_\perp=163.2^\circ=\psi+90^\circ$, $\chi_\perp=-21.9^\circ$ (izquierda), $\mathbf S_\perp=(1,\,0.6,\,-0.4,\,-0.693)$.}

\codigo{t3_p18_elipse_poincare.py}
\begin{figure}[H]\centering
\includegraphics[width=0.95\linewidth]{t3_p18_elipse_poincare.png}
\caption{Problema 18. Izquierda: elipses de $\J$ (derecha, giro horario) y de $\J_\perp$ (izquierda,
giro antihorario) en $z=0$, con la onda viniendo hacia el lector; las líneas punteadas son los ejes
mayores, a $73.2^\circ$ y $163.2^\circ$. Derecha: los dos estados son puntos antípodas en la esfera
de Poincaré.}
\label{fig:t3_p18}
\end{figure}

\idea{(i) $\tan2\psi$ tiene dos soluciones; el cuadrante lo dan los signos de $\mathsf S_1$ (que va
con $\cos2\psi$) y $\mathsf S_2$ (que va con $\sin2\psi$). (ii) Para el sentido de giro, en Saleh basta el signo de
$\sin\varphi$ ($\mathsf S_3>0$: derecha, hemisferio norte). Otros textos (ingeniería, Boyd con
$e^{-i\omega t}$) usan el nombre contrario: antes de decir ``derecha'' hay que fijar la convención. (iii) Ortogonal significa antípoda
en Poincaré: ejes perpendiculares y giro opuesto. El estado ortogonal de $(A_x,A_y)$ es
$(-A_y^*,A_x^*)$; no olvide el conjugado.}

%---------------------------------------------------------------------
\sol{19}{Retardadores y rotadores}{Saleh, §6.1B, ecs.~6.1-19 a 6.1-26, pp.~206--209; Probs.~6.1-6, 6.1-7, p.~240; Ej.~6.1-4, p.~209; ec.~6.6-1, p.~237.}
\label{sol:19}

\inc{a} Comparando (6.1-20) con (6.1-22) se ve que el rotador de ángulo $\theta$ es
$\T_{\rm rot}(\theta)=\R(-\theta)$. Las matrices de rotación en el plano se componen sumando
ángulos:
\[
\R(a)\R(b)=\begin{bmatrix}\cos a\cos b-\sin a\sin b & \cos a\sin b+\sin a\cos b\\
-\sin a\cos b-\cos a\sin b & -\sin a\sin b+\cos a\cos b\end{bmatrix}
=\begin{bmatrix}\cos(a+b)&\sin(a+b)\\-\sin(a+b)&\cos(a+b)\end{bmatrix}=\R(a+b),
\]
y por eso conmutan. En unos ejes girados un ángulo $\phi$, la matriz del rotador es (ec.~6.1-23)
\[
\T'=\R(\phi)\,\T_{\rm rot}(\theta)\,\R(-\phi)=\R(\phi)\R(-\theta)\R(-\phi)=\R(\phi-\theta-\phi)=\R(-\theta)=\T_{\rm rot}(\theta).
\]
\resultado{La matriz del rotador no depende de la orientación de los ejes: un rotador no tiene
``eje'' privilegiado.}

\inc{b} Con el eje rápido en $x$, el retardador $\lambda/2$ ($\Gamma=\pi$) es
$\T=\mathrm{diag}(1,e^{-j\pi})=\mathrm{diag}(1,-1)$. Si la luz es LP a un ángulo $\theta$ del eje rápido:
\[
\begin{bmatrix}1&0\\0&-1\end{bmatrix}\begin{bmatrix}\cos\theta\\ \sin\theta\end{bmatrix}
=\begin{bmatrix}\cos\theta\\ -\sin\theta\end{bmatrix}=\begin{bmatrix}\cos(-\theta)\\ \sin(-\theta)\end{bmatrix}:
\]
sale LP a $-\theta$, es decir, girada $2\theta$ (hacia el otro lado del eje rápido).
\emph{Por qué no es un rotador:}
\begin{itemize}
\item El ángulo girado, $2\theta$, depende de la polarización de entrada; un rotador gira
\emph{toda} LP el mismo ángulo. La lámina $\lambda/2$ hace un \emph{reflejo} respecto al eje rápido,
no una rotación: $\det\T=-1$, mientras que $\det\T_{\rm rot}=+1$.
\item Con luz circular: $\mathrm{diag}(1,-1)\,(1,j)^T=(1,-j)^T$. La RCP sale LCP: la lámina invierte
el sentido de giro. Un rotador deja la circular intacta (salvo una fase):
$\T_{\rm rot}(\theta)(1,j)^T=(\cos\theta-j\sin\theta,\ \sin\theta+j\cos\theta)^T=e^{-j\theta}(1,j)^T$.
\end{itemize}
\resultado{La lámina $\lambda/2$ refleja la polarización respecto a su eje rápido ($\theta\to-\theta$, R$\leftrightarrow$L); el rotador gira todo el mismo ángulo y conserva la circularidad.}

\inc{c} Los modos normales son los autovectores, $\T\J=\mu\J$ (ec.~6.1-26). Los autovalores
salen de $\mu^2-(\mathrm{tr}\,\T)\mu+\det\T=0$.

\emph{Retardador $\Gamma$ con eje rápido a $30^\circ$.} De \eqref{eq:t3_retgen} con $c=\sqrt3/2$, $s=1/2$:
\[
\T=\begin{bmatrix}\frac34+\frac14e^{-j\Gamma} & \frac{\sqrt3}{4}(1-e^{-j\Gamma})\\[2pt]
\frac{\sqrt3}{4}(1-e^{-j\Gamma}) & \frac14+\frac34e^{-j\Gamma}\end{bmatrix},\qquad
\mathrm{tr}\,\T=1+e^{-j\Gamma},\quad \det\T=e^{-j\Gamma}.
\]
Entonces $\mu^2-(1+e^{-j\Gamma})\mu+e^{-j\Gamma}=(\mu-1)(\mu-e^{-j\Gamma})=0$, y $\mu_1=1$,
$\mu_2=e^{-j\Gamma}$. Los autovectores salen sin cuentas: $\J_1=(\cos30^\circ,\sin30^\circ)=\R(-30^\circ)\hat e_x$, y
$\T\J_1=\R(-30^\circ)\,\mathrm{diag}(1,e^{-j\Gamma})\,\R(30^\circ)\R(-30^\circ)\hat e_x=\R(-30^\circ)\hat e_x=\J_1$.
Del mismo modo, $\J_2=\R(-30^\circ)\hat e_y=(-\sin30^\circ,\cos30^\circ)$, LP a $120^\circ$, con
$\mu_2=e^{-j\Gamma}$.

\emph{Rotador de ángulo $\theta$.} $\mathrm{tr}=2\cos\theta$ y $\det=1$, así que
$\mu^2-2\mu\cos\theta+1=0$ y $\mu=\cos\theta\pm j\sin\theta=e^{\pm j\theta}$.
Para $\mu=e^{-j\theta}$, la primera fila da $(\cos\theta-e^{-j\theta})A_x-\sin\theta A_y=0$, es decir,
$j\sin\theta\,A_x=\sin\theta\,A_y$, y por tanto $A_y=jA_x$: \textbf{RCP}. Para
$\mu=e^{+j\theta}$: $A_y=-jA_x$, \textbf{LCP}. (Coincide con lo visto en (b).)

\emph{Polarizador a $45^\circ$.} De (6.1-25), $\T=\frac12\left[\begin{smallmatrix}1&1\\1&1\end{smallmatrix}\right]$:
$\mathrm{tr}=1$ y $\det=0$, así que $\mu(\mu-1)=0$. Para $\mu=1$, $\J=\frac{1}{\sqrt2}(1,1)$ (LP a $45^\circ$,
pasa entera); para $\mu=0$, $\J=\frac{1}{\sqrt2}(-1,1)$ (LP a $135^\circ$, se bloquea).
\resultado{Retardador a $30^\circ$: LP a $30^\circ$ ($\mu=1$) y a $120^\circ$ ($\mu=e^{-j\Gamma}$).
Rotador: RCP ($\mu=e^{-j\theta}$) y LCP ($\mu=e^{+j\theta}$). Polarizador a $45^\circ$: LP a $45^\circ$ ($\mu=1$) y a $135^\circ$ ($\mu=0$).}

\inc{d} El retardador está a $45^\circ$, y por \eqref{eq:t3_retgen} con $c=s=1/\sqrt2$:
\[
\T_{\rm ret}=\frac12\begin{bmatrix}1+e^{-j\Gamma} & 1-e^{-j\Gamma}\\ 1-e^{-j\Gamma} & 1+e^{-j\Gamma}\end{bmatrix},
\qquad \T_{\rm pol\,y}=\begin{bmatrix}0&0\\0&1\end{bmatrix}.
\]
La cascada es $\T=\T_{\rm pol\,y}\T_{\rm ret}$ (el primer elemento atravesado va a la derecha, Saleh, p.~207).
Con entrada $(1,0)^T$:
\[
\J_{\rm sal}=\T_{\rm pol\,y}\,\frac12\begin{bmatrix}1+e^{-j\Gamma}\\ 1-e^{-j\Gamma}\end{bmatrix}
=\begin{bmatrix}0\\ \tfrac12(1-e^{-j\Gamma})\end{bmatrix}
=\begin{bmatrix}0\\ \tfrac12e^{-j\Gamma/2}(e^{j\Gamma/2}-e^{-j\Gamma/2})\end{bmatrix}
=\begin{bmatrix}0\\ j\,e^{-j\Gamma/2}\sin(\Gamma/2)\end{bmatrix}.
\]
La transmitancia de potencia es $\mathcal T=|J_{{\rm sal},y}|^2/|J_{\rm ent}|^2=\sin^2(\Gamma/2)$.
\resultado{$\mathcal T=\sin^2(\Gamma/2)$, idéntica a (6.6-1) (Saleh, p.~237). Es el mismo montaje:
la luz LP en $x$ hace las veces del primer polarizador, el polarizador en $y$ está cruzado con él, y el
retardador está a $45^\circ$ de ambos.}

\idea{Para modos normales y autovalores use traza y determinante: $\mu_1+\mu_2=\mathrm{tr}\,\T$ y
$\mu_1\mu_2=\det\T$. Retardador: modos LP en sus ejes. Rotador: modos circulares. Polarizador:
$\mu=1,0$. El orden de la cascada es de derecha a izquierda. La lámina $\lambda/2$ ``parece'' un rotador
para \emph{una} LP dada, pero es un reflejo ($\det=-1$) e invierte R$\leftrightarrow$L.}

%---------------------------------------------------------------------
\sol{20}{Reflexión de luz circular y pantalla antirreflejo}{Saleh, Probs.~6.1-9 y 6.1-10, p.~241; §6.2, Fig.~6.2-1, pp.~209--211; §6.6D, p.~238.}
\label{sol:20}

\inc{a} \emph{La trampa de los ejes.} El sentido ``derecho/izquierdo'' se define respecto a la
dirección de propagación: el observador mira la onda que \emph{viene hacia él} (Saleh, p.~201). Para
la onda reflejada esa dirección es la contraria, y hay que describirla en un sistema dextrógiro
$(x',y',z')$ con $z'=-z$, por ejemplo $x'=x$, $y'=-y$ (se verifica que $\hat x'\times\hat y'=\hat x\times(-\hat y)=-\hat z=\hat z'$).

\emph{Cálculo.} Sea una RCP que incide normalmente desde $z<0$ sobre un espejo perfecto en $z=0$:
$\mathbf E_1=A(\hat x+j\hat y)e^{-jkz}$. En la superficie de un conductor perfecto el campo eléctrico tangencial
total se anula (Saleh, Ej.~6.2-2: $\mathcal R\to1$ si $\sigma\to\infty$; con
(6.2-8), (6.2-9) y $n_2\to\infty$, $\mathsf r_x=\mathsf r_y=-1$ en incidencia normal). Entonces, en
ejes del laboratorio,
\[
\mathbf E_3=-A(\hat x+j\hat y)e^{+jkz}.
\]
Escrita en los ejes $(x',y')$ propios de la onda reflejada, $E_{3x'}=E_{3x}=-A$ y
$E_{3y'}=-E_{3y}=+jA$, es decir,
\[
\J_3'=-A\begin{bmatrix}1\\-j\end{bmatrix}\ \Rightarrow\ \varphi'=-\pi/2:\ \text{LCP}.
\]
\emph{Por qué.} En el laboratorio el campo en $z=0$ sigue girando en el mismo sentido: con
$\Ec_3=-A(\cos\omega t,\,-\sin\omega t)$, el campo pasa de $(-A,0)$ a $(0,A)$, horario visto desde $+z$.
Pero ahora la onda viaja hacia $-z$, y quien la recibe la mira desde $-z$, donde ese mismo giro es
antihorario. El giro en el espacio se conserva y la dirección de propagación se invierte, así que
la helicidad cambia. Lo mismo pasa en incidencia normal sobre cualquier dieléctrico, porque
$\mathsf r_x=\mathsf r_y$ (Saleh, p.~213).

\emph{Nota sobre los ejes de Saleh.} En incidencia normal, la Fig.~6.2-1 usa para la onda reflejada
los mismos $x$ e $y$ del laboratorio, y por eso obtiene $\mathsf r_x=\mathsf r_y=(n_1-n_2)/(n_1+n_2)$.
Ese sistema es \emph{levógiro} respecto a $\kv_3$. Si se usa $\J_3=\mathbf r\J_1$ tal cual, el vector
$(1,j)$ parece seguir siendo ``RCP'', y eso es falso. Esta misma elección de ejes explica el $\pi$ extra
de $\mathsf r_y$ frente a la convención de Born \& Wolf en el Problema~21(c).
\resultado{Una RCP se refleja como LCP (y viceversa): el sentido de giro en el espacio se
conserva, la dirección de propagación se invierte.}

\inc{b} \emph{Montaje.} De afuera hacia la ventana: polarizador lineal en $x$ y lámina $\lambda/4$
con eje rápido a $45^\circ$. Todo se escribe en ejes fijos del laboratorio, los mismos para la ida y
para la vuelta. En esos ejes, un elemento \emph{recíproco} tiene como matriz de vuelta la
\emph{transpuesta} de la de ida (si la onda vuelve, recorre los cambios al revés, Saleh §6.6D,
p.~238). Polarizador y retardador tienen matrices simétricas, de modo que la matriz de vuelta es la misma.

Por \eqref{eq:t3_retgen} con $\theta=45^\circ$ y $e=e^{-j\pi/2}=-j$:
\[
\T_Q=\frac12\begin{bmatrix}1-j&1+j\\1+j&1-j\end{bmatrix},\qquad
\T_Q^2=\T_{\rm ret}(\pi,45^\circ)=\begin{bmatrix}0&1\\1&0\end{bmatrix}.
\]
\emph{Ida.} La luz de fondo (no polarizada) atraviesa el polarizador, $\J=(1,0)^T$ con la mitad de la
potencia, y luego la lámina:
$\T_Q(1,0)^T=\frac12(1-j,\,1+j)^T=\frac{1-j}{2}(1,\,j)^T$, pues $(1+j)/(1-j)=j$. Sale RCP.

\emph{Reflexión en la ventana} (casi normal): $\mathsf r_x=\mathsf r_y=\mathsf r\approx-0.2$, igual para las dos
componentes del laboratorio. Por (a), la onda vuelve como LCP respecto a su propia dirección.

\emph{Vuelta.} La lámina convierte $(1,j)$ en
$\T_Q(1,j)^T=\frac12\big((1-j)+j(1+j),\ (1+j)+j(1-j)\big)^T=(0,\ 1+j)^T$, que es LP en $y$, y el polarizador en $x$ la bloquea.
De un solo golpe: $\T_{\rm total}=\mathsf r\,\T_{\rm pol\,x}\T_Q^2\T_{\rm pol\,x}
=\mathsf r\left[\begin{smallmatrix}1&0\\0&0\end{smallmatrix}\right]\left[\begin{smallmatrix}0&1\\1&0\end{smallmatrix}\right]\left[\begin{smallmatrix}1&0\\0&0\end{smallmatrix}\right]=0$.
La ida y la vuelta por la $\lambda/4$ equivalen a una $\lambda/2$ a $45^\circ$, que lleva $x$ a $y$.
La luz del objeto autoluminoso, en cambio, cruza la ventana una sola vez hacia el observador y pasa
(con la mitad de su potencia si es no polarizada).

\emph{¿Es un aislador óptico?} Solo en un sentido restringido. Bloquea la luz que vuelve después de
reflejarse en un reflector que conserva el vector de Jones del laboratorio (incidencia casi normal, espejos).
Pero es un dispositivo recíproco. Una onda que llegue desde la ventana en el estado de laboratorio
$(1,-j)$ pasa entera: $\T_Q(1,-j)^T=(1-j,\,0)^T$, que el polarizador transmite. Además, en incidencia oblicua
$\mathsf r_{TE}\neq\mathsf r_{TM}$, la onda reflejada ya no es circular y el bloqueo es parcial. El aislador de
Faraday (Saleh, Fig.~6.6-5, p.~238) bloquea \emph{cualquier} onda que regrese, sin importar su
polarización ni cómo se generó, porque el rotador de Faraday no es recíproco (Problema~26).
\resultado{Polarizador + $\lambda/4$ a $45^\circ$: ida y vuelta equivalen a una $\lambda/2$ a $45^\circ$
($x\to y$), y el reflejo se bloquea. No es un aislador verdadero: es recíproco y solo funciona porque la
reflexión invierte la helicidad. El de Faraday bloquea todo lo que regresa.}

\idea{Para ida y vuelta, trabaje en ejes \emph{fijos} del laboratorio. Elemento recíproco: matriz
de vuelta $=\T^T$. Faraday: la misma $\T$. Para decidir ``derecha/izquierda'' después de una
reflexión, cambie a un sistema dextrógiro de la onda reflejada (invierta un eje).}

%---------------------------------------------------------------------
\sol{21}{Fresnel: reflectancias y desfases}{Saleh, §6.2, ecs.~6.2-1 a 6.2-17, pp.~209--214; Probs.~6.2-3, 6.2-4, 6.2-6, p.~241.}
\label{sol:21}

\inc{a} \emph{Geometría.} Frontera $z=0$, medio 1 en $z<0$, plano de incidencia $y$--$z$. La onda
TE tiene $\mathbf E\parallel\hat x$ (perpendicular al plano de incidencia). Los vectores unitarios de
propagación son
\[
\hat k_1=\sin\theta_1\hat y+\cos\theta_1\hat z,\quad
\hat k_3=\sin\theta_1\hat y-\cos\theta_1\hat z,\quad
\hat k_2=\sin\theta_2\hat y+\cos\theta_2\hat z .
\]
Los campos son $\mathbf E_i=A_{ix}\hat x\,e^{-j\kv_i\cdot\rv}$, y para ondas planas
$\Hv=\hat k\times\mathbf E/\eta$ (Saleh, §5.4, ec.~5.4-5). El ajuste de fase en $z=0$, para todo $y$, exige
$k_1\sin\theta_1=k_2\sin\theta_2$ (Snell, 6.2-1) y $\theta_3=\theta_1$.

\emph{Componentes tangenciales} ($x$ para $\mathbf E$, $y$ para $\Hv$). Como
$\hat k\times\hat x=(\sin\theta\,\hat y+\cos\theta\,\hat z)\times\hat x=-\sin\theta\,\hat z+\cos\theta\,\hat y$,
\[
H_{1y}=\frac{\cos\theta_1}{\eta_1}A_{1x},\qquad H_{3y}=-\frac{\cos\theta_1}{\eta_1}A_{3x},\qquad
H_{2y}=\frac{\cos\theta_2}{\eta_2}A_{2x}.
\]
\emph{Continuidad} en $z=0$ (Saleh, p.~210):
\begin{align*}
E_t:&\quad A_{1x}+A_{3x}=A_{2x},\\
H_t:&\quad \frac{\cos\theta_1}{\eta_1}(A_{1x}-A_{3x})=\frac{\cos\theta_2}{\eta_2}A_{2x}.
\end{align*}
Se sustituye la primera en la segunda:
$\frac{\cos\theta_1}{\eta_1}(A_{1x}-A_{3x})=\frac{\cos\theta_2}{\eta_2}(A_{1x}+A_{3x})$, y se agrupa:
\[
A_{1x}\Big(\frac{\cos\theta_1}{\eta_1}-\frac{\cos\theta_2}{\eta_2}\Big)=A_{3x}\Big(\frac{\cos\theta_1}{\eta_1}+\frac{\cos\theta_2}{\eta_2}\Big)
\;\Rightarrow\;
\mathsf r_x=\frac{A_{3x}}{A_{1x}}=\frac{\cos\theta_1/\eta_1-\cos\theta_2/\eta_2}{\cos\theta_1/\eta_1+\cos\theta_2/\eta_2}.
\]
Multiplicando numerador y denominador por $\eta_1\eta_2\sec\theta_1\sec\theta_2$:
\[
\boxed{\mathsf r_x=\frac{\eta_2\sec\theta_2-\eta_1\sec\theta_1}{\eta_2\sec\theta_2+\eta_1\sec\theta_1}},\qquad
\mathsf t_x=\frac{A_{2x}}{A_{1x}}=1+\mathsf r_x ,
\]
que es (6.2-6). Con $\eta=\eta_0/n$ resulta
$\mathsf r_x=(n_1\cos\theta_1-n_2\cos\theta_2)/(n_1\cos\theta_1+n_2\cos\theta_2)$, la ec.~(6.2-8).
(Si la luz es un haz y no una onda plana, se descompone en ondas planas, se aplica $\mathsf r(\theta)$ a cada una
y se resuma; de ahí sale el Goos--Hänchen del Problema~22.)
\resultado{$\mathsf r_x=(\eta_2\sec\theta_2-\eta_1\sec\theta_1)/(\eta_2\sec\theta_2+\eta_1\sec\theta_1)$, $\mathsf t_x=1+\mathsf r_x$ (6.2-6).}

\inc{b} $n_1=1$, $n_2=1.5$, $\theta_1=45^\circ$. Por Snell, $\sin\theta_2=0.7071/1.5=0.4714$ y
$\cos\theta_2=0.8819$.
\begin{align*}
\mathsf r_{TE}&=\frac{1\cdot0.7071-1.5\cdot0.8819}{0.7071+1.3229}=\frac{-0.6158}{2.0300}=-0.3033,
&\mathcal R_{TE}&=|\mathsf r_{TE}|^2=0.0920,\\
\mathsf r_{TM}&=\frac{1/0.7071-1.5/0.8819}{1/0.7071+1.5/0.8819}=\frac{1.4142-1.7008}{1.4142+1.7008}=-0.0920,
&\mathcal R_{TM}&=0.0085 .
\end{align*}
Luz no polarizada, con TE y TM de igual intensidad: $\bar{\mathcal R}=\frac12(0.0920+0.0085)=0.0502$.
Comprobación: a $45^\circ$ se cumple la identidad $\mathsf r_{TM}=\mathsf r_{TE}^2$ ($0.3033^2=0.0920$).
Orden de magnitud: en incidencia normal $\mathcal R=0.04$ (6.2-15); a $45^\circ$ la TE sube y la TM baja,
porque se acerca a Brewster ($56.3^\circ$).
\resultado{$\mathcal R_{TE}=0.0920$, $\mathcal R_{TM}=0.0085$, $\bar{\mathcal R}=0.0502$ (5.0\,\%).}

\inc{c} \emph{Datos.} $\theta_c=\sin^{-1}(1/1.5)=41.81^\circ$ y $\theta_1=1.2\theta_c=50.17^\circ$. Con la rama de
Saleh (p.~212, nota al pie),
\[
\cos\theta_2=-j\sqrt{\frac{\sin^2\theta_1}{\sin^2\theta_c}-1}=-j\sqrt{\Big(\frac{0.7682}{0.6667}\Big)^2-1}=-0.5719\,j .
\]
(Esa rama hace que la onda transmitida, $\propto e^{-jk_2z\cos\theta_2}=e^{-k_2(0.5719)z}$, \emph{decaiga}.)

\emph{TE} (6.2-8), con $a=n_1\cos\theta_1=1.5\cdot0.6405=0.9607$ y $b=n_2\cdot0.5719=0.5719$:
\[
\mathsf r_x=\frac{a-n_2\cos\theta_2}{a+n_2\cos\theta_2}=\frac{a+jb}{a-jb}\ \Rightarrow\ |\mathsf r_x|=1,\quad
\varphi_{TE}=2\tan^{-1}\frac{b}{a}=2\tan^{-1}(0.5952)=+61.53^\circ .
\]
Coincide con (6.2-11), $\tan(\varphi_x/2)=\sqrt{\cos^2\theta_c/\cos^2\theta_1-1}\equiv X=0.5952$, porque
\[
b/a=\sqrt{\sin^2\theta_1-\sin^2\theta_c}\,/\cos\theta_1=\sqrt{\cos^2\theta_c-\cos^2\theta_1}\,/\cos\theta_1=X .
\]

\emph{TM} (6.2-9). $\sec\theta_2=1/(-0.5719j)=+1.7486\,j$ y $n_1\sec\theta_1=1.5/0.6405=2.3420$:
\[
\mathsf r_y=\frac{2.3420-1.7486\,j}{2.3420+1.7486\,j}\ \Rightarrow\ |\mathsf r_y|=1,\quad
\varphi_{TM}=-2\tan^{-1}\frac{1.7486}{2.3420}=-2\tan^{-1}(0.7466)=-73.50^\circ .
\]
En forma cerrada: el cociente es $\dfrac{n_2|\sec\theta_2|}{n_1\sec\theta_1}=\dfrac{\sin^2\theta_c\cos\theta_1}{\sqrt{\cos^2\theta_c-\cos^2\theta_1}}=\dfrac{\sin^2\theta_c}{X}$, de modo que
\begin{equation}
\tan\frac{\varphi_{TM}}{2}=-\frac{\sin^2\theta_c}{\sqrt{\cos^2\theta_c/\cos^2\theta_1-1}}\qquad\text{(forma que sale de 6.2-9)}.
\label{eq:t3_tm}
\end{equation}
\emph{Retardo:} $\varphi_{TE}-\varphi_{TM}=61.53^\circ-(-73.50^\circ)=135.0^\circ$.

\emph{La fórmula impresa (6.2-13).} Saleh imprime (p.~213)
$\tan(\varphi_y/2)=-\dfrac{1}{\sin^2\theta_c}\sqrt{\cos^2\theta_c/\cos^2\theta_1-1}$, que es el
\emph{inverso} del cociente de \eqref{eq:t3_tm}. Las dos expresiones no son equivalentes:
\begin{itemize}
\item En $\theta_1=\theta_c$, $X=0$: la impresa da $\tan(\varphi_y/2)=0$, es decir, $\varphi_y=0$. Pero
directamente, $\cos\theta_2=0$ hace $\sec\theta_2\to\infty$ y $\mathsf r_y\to-1$, es decir, $\varphi_y=\pm\pi$.
Eso es lo que muestra la Fig.~6.2-5 (p.~213): la fase vale $\pi$ en $\theta_c$ y baja a 0 en
$90^\circ$. También lo dice el texto de la misma página: $\mathsf r_y$ es negativo y llega a $|\mathsf r_y|=1$
en $\theta_c$.
\item En $1.2\theta_c$ la impresa da $\varphi_y=2\tan^{-1}(-X/\sin^2\theta_c)=-106.5^\circ$; la directa,
$-73.5^\circ$. Llamando $P$ a la fase impresa y $D$ a la directa, $\tan(P/2)\tan(D/2)=1$, y por eso $P=-\pi-D$.
\item La magnitud de la curva de la figura en $50.2^\circ$ es $\approx0.41\pi=73^\circ$, que coincide con
$|\varphi_{TM}|$ directo. La figura dibuja la fase como un retraso positivo de $\pi$ a 0; la curva
directa vale $-\pi\equiv\pi$ en $\theta_c$ y 0 en $90^\circ$, así que coincide con la figura en magnitud y en los extremos.
\end{itemize}
Conclusión: la (6.2-13) impresa \emph{no} es consistente con (6.2-9) ni con la Fig.~6.2-5. Para el
examen, calcule las fases directamente con Fresnel (6.2-8)--(6.2-9) y la rama
$\cos\theta_2=-j\sqrt{\cdots}$, o use \eqref{eq:t3_tm}.

\emph{Convención de Born \& Wolf.} Saleh orienta los ejes $y$ de las ondas incidente y reflejada
de modo que sus componentes paralelas a la frontera coincidan (Fig.~6.2-1). Por eso en incidencia
normal $\mathsf r_x=\mathsf r_y$ (ver Problema~20(a)). Born \& Wolf usan ejes dextrógiros para cada
onda, y su $r_{TM}$ tiene el signo opuesto: $r_{TM}^{BW}=-\mathsf r_y$, es decir,
$\varphi_{TM}^{BW}=\varphi_{TM}+\pi=106.5^\circ$. Entonces
\[
\delta\equiv\varphi_{TE}-\varphi_{TM}^{BW}=61.53^\circ-106.50^\circ=-45.0^\circ,\qquad |\delta|=45.0^\circ .
\]
Con las expresiones cerradas, $\delta/2=\tan^{-1}X+\tan^{-1}(\sin^2\theta_c/X)-\pi/2=\tan^{-1}X-\tan^{-1}(X/\sin^2\theta_c)$, y
\[
\tan\frac{\delta}{2}=\frac{X-X/s^2}{1+X^2/s^2}=-\frac{X\cos^2\theta_c}{s^2+X^2}
=-\frac{\cos\theta_1\sqrt{\sin^2\theta_1-\sin^2\theta_c}}{\sin^2\theta_1}\qquad(s\equiv\sin\theta_c),
\]
donde se usó $s^2+X^2=\cos^2\theta_c\sin^2\theta_1/\cos^2\theta_1$ y $X\cos\theta_1=\sqrt{\sin^2\theta_1-s^2}$.
Es la fórmula clásica del retardo en RTI. Numéricamente:
$\cos\theta_1\sqrt{\sin^2\theta_1-s^2}/\sin^2\theta_1=0.6405\cdot0.3817/0.5901=0.4143=\tan22.5^\circ$, así que $|\delta|=45.0^\circ$.
La diferencia entre $135^\circ$ y $45^\circ$ es solo el $\pi$ de la convención (con el signo); la física es la misma.
\resultado{$\theta_c=41.8^\circ$, $\theta_1=50.2^\circ$: $\varphi_{TE}=+61.5^\circ$, $\varphi_{TM}=-73.5^\circ$,
$\varphi_{TE}-\varphi_{TM}=135.0^\circ$ (Saleh), es decir, $45.0^\circ$ en la convención de Born \& Wolf. La (6.2-13) impresa da
$\varphi_y=0$ en $\theta_c$, contra el $\pi$ de la figura y de (6.2-9): está invertida; use \eqref{eq:t3_tm}.}

\inc{d} \textbf{[N]} \emph{Rombo de Fresnel.} Se calcula $\delta(\theta)=\varphi_{TE}-\varphi_{TM}-\pi$ con Fresnel
directo para $n=1.51$ y se compara con la fórmula cerrada anterior; coinciden con un error de $10^{-11}$
grados. El retardo vale 0 en $\theta_c=41.47^\circ$ y en $90^\circ$, y tiene un máximo de
$45.94^\circ$ en $51.3^\circ$. Como el máximo apenas pasa de $45^\circ$, hay dos ángulos con
$\delta=45^\circ$:
\[
\theta=48.6^\circ\quad\text{y}\quad\theta=54.6^\circ .
\]
Con dos reflexiones a uno de esos ángulos, el retardo total es $90^\circ$: un retardador $\lambda/4$. En la
convención de Saleh cada reflexión da $135^\circ$ y dos dan $270^\circ\equiv-90^\circ$. También es
cuarto de onda, porque el $\pi$ extra de cada reflexión suma $2\pi$. El rombo comercial usa $54.6^\circ$:
cerca de él, $d\delta/d\theta$ es pequeña y el ángulo exacto importa menos. Es \emph{acromático}
porque $\delta$ depende solo de $n$, que varía poco con $\lambda$, y no de un espesor dividido entre
$\lambda$, como en una lámina birrefringente ($\Gamma=2\pi\Delta n\,d/\lambda_0$).

\codigo{t3_p21_fresnel.py}
\begin{figure}[H]\centering
\includegraphics[width=0.75\linewidth]{t3_p21_rombo_fresnel.png}
\caption{Problema 21(d). Retardo por reflexión total interna (convención de Born \& Wolf) contra el ángulo
interno, para $n=1.51$. La recta a $45^\circ$ corta la curva en $48.6^\circ$ y $54.6^\circ$.}
\end{figure}
\resultado{$n=1.51$: $\delta=45^\circ$ por reflexión en $48.6^\circ$ y $54.6^\circ$ (máximo de $45.9^\circ$ en $51.3^\circ$).}

\inc{e} \textbf{[N]} Con (6.2-8)--(6.2-9) y $\mathcal R=|\mathsf r|^2$ (6.2-14) para $n_2=3.48$ (Si) y $1.99$
(Si$_3$N$_4$):
\[
\theta_B=\tan^{-1}n_2:\quad 74.0^\circ\ (\text{Si}),\qquad 63.3^\circ\ (\text{Si}_3\text{N}_4);\qquad
\mathcal R(0)=\Big(\frac{n_2-1}{n_2+1}\Big)^2=0.306\ (\text{Si}),\ 0.110\ (\text{Si}_3\text{N}_4).
\]
La TE crece de forma monótona hasta 1 en incidencia rasante. La TM baja hasta anularse en $\theta_B$ y luego sube
bruscamente. El Si refleja casi un tercio de la luz en incidencia normal; por eso las facetas de
chips de Si necesitan antirreflejo, o la luz se acopla con rejillas.
\begin{figure}[H]\centering
\includegraphics[width=0.75\linewidth]{t3_p21_reflectancias.png}
\caption{Problema 21(e). Reflectancias de potencia TE (continua) y TM (discontinua) contra el ángulo de incidencia
para aire--Si ($n=3.48$) y aire--Si$_3$N$_4$ ($n=1.99$) a $\lambda_0=\SI{1.55}{\micro\metre}$.}
\end{figure}
\resultado{$\theta_B=74.0^\circ$ (Si) y $63.3^\circ$ (Si$_3$N$_4$); $\mathcal R(0)=30.6$\,\% y $11.0$\,\%.}

\idea{(i) En RTI $|\mathsf r|=1$ y solo importa la fase; tome $\cos\theta_2=-j\sqrt{\cdots}$ (onda
evanescente que decae). (ii) La fase TM de Saleh tiene un $\pi$ de diferencia con la de Born \& Wolf por la
orientación de los ejes $y$; un retardo de ``$135^\circ$'' en Saleh es ``$45^\circ$'' en B\&W. (iii) La
(6.2-13) impresa está invertida: verifique siempre una fórmula en un punto conocido ($\theta_c$:
$\mathsf r_y=-1$). (iv) Luz no polarizada: promedie las \emph{reflectancias}, no las amplitudes.}

%---------------------------------------------------------------------
\sol{22}{Goos--Hänchen y medio absorbente}{Saleh, Probs.~6.2-7 y 6.2-8, p.~241; ec.~6.2-11, p.~212; §5.5A, ecs.~5.5-2 a 5.5-5, pp.~171--172.}
\label{sol:22}

\inc{a} \emph{Coeficiente $\xi$.} La fase TE en RTI es $\varphi(\theta)=2\tan^{-1}X$, con
$X=\sqrt{\cos^2\theta_c/\cos^2\theta-1}=\sqrt{\sin^2\theta-\sin^2\theta_c}/\cos\theta$ (6.2-11). Se deriva:
\[
\frac{d\varphi}{d\theta}=\frac{2}{1+X^2}\frac{dX}{d\theta},\qquad 1+X^2=\frac{\cos^2\theta_c}{\cos^2\theta}.
\]
Con $\sigma\equiv\sqrt{\sin^2\theta-\sin^2\theta_c}$ y la regla del cociente:
\[
\frac{dX}{d\theta}=\frac{\dfrac{\sin\theta\cos\theta}{\sigma}\cos\theta+\sigma\sin\theta}{\cos^2\theta}
=\frac{\sin\theta\,(\cos^2\theta+\sin^2\theta-\sin^2\theta_c)}{\sigma\cos^2\theta}
=\frac{\sin\theta\cos^2\theta_c}{\sigma\cos^2\theta}.
\]
Por tanto
\[
\boxed{\xi=\frac{d\varphi}{d\theta}=\frac{2\cos^2\theta}{\cos^2\theta_c}\cdot\frac{\sin\theta\cos^2\theta_c}{\sigma\cos^2\theta}
=\frac{2\sin\theta}{\sqrt{\sin^2\theta-\sin^2\theta_c}}}
\]
(adimensional, rad/rad). Diverge en $\theta_c$ y tiende a $2/\cos\theta_c$ en incidencia rasante. Para vidrio--aire a
$1.2\theta_c$: $\xi=2(0.7682)/\sqrt{0.5901-0.4444}=4.03$ (la derivada numérica del script da lo mismo).

\emph{Desplazamiento lateral.} Dos ondas planas incidentes a $\theta$ y $\theta+d\theta$, con
amplitudes iguales y fase espacial $e^{-jk_1y\sin\theta}$ a lo largo de la frontera, interfieren con un patrón
$\propto1+\cos(\Delta k_y\,y)$, donde $\Delta k_y=k_1\cos\theta\,d\theta$ (franjas de periodo
$2\pi/\Delta k_y$). Las dos reflejadas llevan además las fases $\varphi(\theta)$ y
$\varphi(\theta)+\xi d\theta$, y su patrón es $\propto1+\cos(\Delta k_y\,y-\xi\,d\theta)$. Es el mismo
patrón corrido una distancia $\Delta y=\xi d\theta/\Delta k_y$ sobre la frontera:
\[
\Delta y=\frac{\xi}{k_1\cos\theta}=\frac{2\tan\theta}{\gamma},\qquad \gamma=k_1\sqrt{\sin^2\theta-\sin^2\theta_c},
\]
independiente de $d\theta$. Un haz es una superposición de muchas ondas planas en torno a $\theta$; cada par
se corre lo mismo, y el haz reflejado sale desplazado $\Delta y$ sobre la frontera, es decir,
$D=\Delta y\cos\theta=\xi/k_1$ perpendicular al haz. El desplazamiento es $\propto\xi$: el haz ``entra'' una
distancia del orden de $1/\gamma$ en el medio menos denso antes de volver.

\emph{Penetración.} En el aire,
$\mathbf E_2\propto e^{-jk_2(y\sin\theta_2+z\cos\theta_2)}=e^{-jk_1y\sin\theta}\,e^{-\gamma z}$, con
\[
\gamma=k_0\sqrt{n_1^2\sin^2\theta-n_2^2}=\frac{2\pi}{\SI{633}{nm}}\sqrt{1.5^2(0.7682)^2-1}
=(\SI{9.926e6}{m^{-1}})(0.5719)=\SI{5.68e6}{m^{-1}} .
\]
\resultado{$\xi=2\sin\theta/\sqrt{\sin^2\theta-\sin^2\theta_c}=4.03$; penetración del campo
$1/\gamma=\SI{176}{nm}$ ($\SI{88}{nm}$ para la intensidad). Corrimiento sobre la frontera
$\Delta y=2\tan\theta/\gamma=\SI{422}{nm}$, desplazamiento del haz $D=\xi/k_1=\SI{271}{nm}\approx0.4\lambda_0$:
del orden de la longitud de onda, como se esperaba.}
\codigo{t3_p22_goos_hanchen.py}

\inc{b} \emph{Índice complejo.} En un medio absorbente $k=\beta-\frac{j}{2}\alpha$ (Saleh, ec.~5.5-3), con
$\beta=nk_0$ (5.5-4), y la onda es $Ae^{-\alpha z/2}e^{-j\beta z}$. Se define
\[
N\equiv\frac{k}{k_0}=n-j\frac{\alpha}{2k_0}=n-j\frac{\alpha c_0}{2\omega}\qquad(k_0=\omega/c_0),
\]
que es (5.5-5): $N=\sqrt{\epsilon/\epsilon_0}$.

\emph{Campos.} Incidencia normal desde el vacío ($z<0$) sobre el medio ($z>0$), con $\mathbf E\parallel\hat x$.
De Faraday, $\nabla\times\mathbf E=-j\omega\mu_0\Hv$; para $\mathbf E=\hat xEe^{\mp jkz}$ se obtiene
$\Hv=\pm\hat y\,(k/\omega\mu_0)E$. Con $k/\omega\mu_0=N/\eta_0$:
\begin{align*}
z<0:&\quad E_x=E_1e^{-jk_0z}+E_3e^{+jk_0z},\qquad H_y=\frac{1}{\eta_0}\big(E_1e^{-jk_0z}-E_3e^{+jk_0z}\big),\\
z>0:&\quad E_x=E_2e^{-jNk_0z},\qquad\qquad\quad\ \ H_y=\frac{N}{\eta_0}E_2e^{-jNk_0z}.
\end{align*}
\emph{Continuidad} de $E_x$ y $H_y$ en $z=0$: $E_1+E_3=E_2$ y $E_1-E_3=NE_2$. Restando,
$E_1-E_3=N(E_1+E_3)$, es decir, $E_1(1-N)=E_3(1+N)$:
\[
\mathsf r\equiv\frac{E_3}{E_1}=\frac{1-N}{1+N}=\frac{1-(n-j\alpha c_0/2\omega)}{1+(n-j\alpha c_0/2\omega)},\qquad
\mathcal R=|\mathsf r|^2=\frac{(n-1)^2+(\alpha c_0/2\omega)^2}{(n+1)^2+(\alpha c_0/2\omega)^2}.
\]
\emph{Discrepancia de signo con el enunciado.} El enunciado (copiado del Prob.~6.2-8 de Saleh) pide
$\mathsf r=[(n-j\alpha c_0/2\omega)-1]/[(n-j\alpha c_0/2\omega)+1]=-\,(1-N)/(1+N)$. Para el campo
eléctrico, y con la convención de la propia ec.~6.2-8 en $\theta=0$ ($\mathsf r=(n_1-n_2)/(n_1+n_2)$), el
resultado correcto es $(1-N)/(1+N)$. La expresión del enunciado es el coeficiente de reflexión del
campo \emph{magnético}, $H_3/H_1=(-E_3/\eta_0)/(E_1/\eta_0)=-\mathsf r$, o el de $\mathbf E$ con la
orientación de ejes de Born \& Wolf para TM. La diferencia es un desfase global de $\pi$ y no cambia
$\mathcal R$. Límite de control: si $\alpha=0$, $\mathsf r=(1-n)/(1+n)$, la conocida (6.2-15).
\resultado{$\mathsf r_E=\dfrac{1-(n-j\alpha c_0/2\omega)}{1+(n-j\alpha c_0/2\omega)}$; el del enunciado,
$[(n-j\alpha c_0/2\omega)-1]/[(n-j\alpha c_0/2\omega)+1]$, es $-\mathsf r_E$ (coeficiente de $\Hv$). $\mathcal R$ es el mismo.}

\idea{Goos--Hänchen: un desfase que depende del ángulo, $d\varphi/d\theta\neq0$, corre el haz (igual que
$d\varphi/d\omega$ retrasa un pulso). El tamaño del corrimiento es $\sim\tan\theta/\gamma$, del orden
de la penetración evanescente, unos cientos de nm. Medio absorbente: basta reemplazar $n\to N=n-j\alpha/2k_0$
en las fórmulas de Fresnel (Saleh: $\mathrm{Im}\,k=-\alpha/2$, con signo menos por $e^{j\omega t}$).
Cuidado con el signo global de $\mathsf r$: depende de si se refiere a $\mathbf E$ o a $\Hv$.}

%---------------------------------------------------------------------
\sol{23}{Cristales uniaxiales: retardo, doble refracción y walk-off}{Saleh, §6.3B--E, ecs.~6.3-15, 6.3-21, 6.3-23 a 6.3-25, pp.~219--227; §6.6B, p.~236; Probs.~6.3-1, 6.3-3, 6.3-5, pp.~241--242.}
\label{sol:23}

\inc{a} El retardo de una lámina es $\Gamma=(n_2-n_1)k_0d$ (Saleh, §6.3B, p.~219; §6.6B, p.~236).
Es máximo cuando $\kv\perp$ eje óptico, porque entonces los dos modos tienen índices $n_o$ y
$n(90^\circ)=n_e$, la mayor diferencia posible (6.3-15). Por unidad de longitud:
\[
\frac{\Gamma}{d}=\frac{2\pi(n_e-n_o)}{\lambda_0}=\frac{2\pi(0.009)}{\SI{633e-9}{m}}=\SI{8.93e4}{rad/m}=\SI{89.3}{rad/mm}\ (14.2\ \text{ondas/mm}).
\]
\emph{Cuarto de onda:} $\Gamma=\pi/2+2\pi m$ da
$d=\frac{\lambda_0}{4\Delta n}(1+4m)=\SI{17.58}{\micro\metre}+m(\SI{70.33}{\micro\metre})$, es decir, $17.6$,
$87.9$, $158.3\ \um$\ldots{} También actúa como $\lambda/4$ el retardo $\Gamma=3\pi/2+2\pi m\equiv-\pi/2$
($52.8$, $123.1$, $193.4\ \um$\ldots), con los papeles de los ejes rápido y lento intercambiados (la LP
a $45^\circ$ sale circular del sentido opuesto). En total, $d=(2m+1)\,\SI{17.58}{\micro\metre}$.
Las láminas de \emph{orden cero} ($\SI{17.6}{\micro\metre}$) son frágiles; por eso se fabrican de orden
múltiple, que son más sensibles a $\lambda$ y a la temperatura.
\resultado{$\Gamma/d=\SI{89.3}{rad/mm}$; $\lambda/4$ en $17.6$, $87.9$, $158\ \um$, \ldots{} (y, con ejes intercambiados, $52.8$, $123\ \um$, \ldots).}

\inc{b} \emph{Geometría.} En el plano de incidencia, sea $\hat z'$ la normal hacia el cristal y
$\hat x'$ la tangente. La onda incidente va en $(\sin30^\circ,\cos30^\circ)$ y el eje óptico es
perpendicular a ella, $\hat c=(\cos30^\circ,-\sin30^\circ)$: forma $\theta_a=60^\circ$ con la normal, del
lado opuesto. Una onda refractada a $\theta$ de la normal forma con el eje un ángulo $\psi$ tal que
$\cos\psi=\sin\theta\cos30^\circ-\cos\theta\sin30^\circ=\sin(\theta-30^\circ)$, de donde
$\psi=120^\circ-\theta$. Como $n(\psi)$ depende de $\cos^2\psi$, es lo mismo usar
$\theta_a+\theta=60^\circ+\theta$, que es la forma de (6.3-25).

\emph{Onda ordinaria} (polarización perpendicular al plano de incidencia, TE), por (6.3-24):
$\sin\theta_o=\sin30^\circ/n_o=0.5/1.544=0.32383$, de donde $\theta_o=18.89^\circ$. Para ella $\Sv\parallel\kv$,
así que el rayo va también a $18.89^\circ$.

\emph{Onda extraordinaria} (polarización en el plano de incidencia). Hay que resolver (6.3-25),
$\sin30^\circ=n(60^\circ+\theta_e)\sin\theta_e$, con
$1/n^2(\psi)=\cos^2\psi/n_o^2+\sin^2\psi/n_e^2$ (6.3-15). Por iteración:
\begin{enumerate}
\item Se arranca con $\theta_e\approx\theta_o=18.89^\circ$, de modo que $\psi=78.89^\circ$:
$1/n^2=0.03709/2.38394+0.96291/2.41181=0.414806$, y $n=1.55267$.
\item $\sin\theta_e=0.5/1.55267=0.32202$, es decir, $\theta_e=18.786^\circ$ y $\psi=78.786^\circ$; al recalcular,
$n=1.55266$. Converge.
\end{enumerate}
\emph{Rayo extraordinario.} Es normal a la superficie $\kv$, el elipsoide (6.3-21). Con $\psi$ medido
desde el eje, la dirección del rayo forma con el eje el ángulo $\psi_S$ dado por
$\tan\psi_S=(n_o^2/n_e^2)\tan\psi=0.98844\cdot5.0444=4.9861$, es decir, $\psi_S=78.659^\circ$. El rayo
está a $60^\circ+\theta_S$ del eje, de donde $\theta_S=18.66^\circ$ de la normal. El walk-off es
$\psi-\psi_S=0.13^\circ$: el cuarzo es positivo y el rayo se inclina \emph{hacia} el eje óptico.
\resultado{Onda o: $\kv$ y rayo a $18.89^\circ$ de la normal ($n_o=1.544$). Onda e: $\kv$ a $18.79^\circ$
($n=1.5527$) y rayo a $18.66^\circ$ (walk-off $0.13^\circ$). Los rayos o y e se separan $0.24^\circ$.}

\inc{c} En incidencia normal los dos $\kv$ refractados son normales a la cara (Snell con
$\sin\theta_1=0$), y $\kv$ forma $45^\circ$ con el eje.
\[
\frac{1}{n^2(45^\circ)}=\frac{\cos^245^\circ}{2.20^2}+\frac{\sin^245^\circ}{2.29^2}=\frac12(0.206612+0.190686)=0.198649
\;\Rightarrow\;n(45^\circ)=2.2436 .
\]
\emph{Walk-off.} El rayo e forma con el eje el ángulo $\theta_S$, con
$\tan\theta_S=(n_o^2/n_e^2)\tan45^\circ=4.84/5.2441=0.92294$, es decir, $\theta_S=42.70^\circ$. Entonces
$\rho=45^\circ-42.70^\circ=2.30^\circ$ ($\SI{40.0}{mrad}$), con el rayo inclinado hacia el eje (Saleh, Fig.~6.3-14,
p.~227: ``la incidencia normal produce refracción oblicua''). Forma equivalente:
$\tan\rho=\frac{n^2}{2}\big(\frac{1}{n_o^2}-\frac{1}{n_e^2}\big)\sin2\theta$.

\emph{Desplazamiento.} El rayo o sigue la normal y el e se desvía $\rho$; al salir, los dos vuelven a
ser normales a la cara (sus frentes siguen paralelos a las caras) y quedan separados
\[
\Delta x=L\tan\rho=(\SI{10}{mm})\tan2.30^\circ=\SI{0.40}{mm},
\]
con polarizaciones ortogonales: la lámina es un divisor de haz polarizador (Saleh, Fig.~6.3-15).

\emph{Retardo.} La fase de cada onda la fija $\kv\cdot\rv$ (los frentes son paralelos a las caras), así que
\[
\Gamma=(n(45^\circ)-n_o)k_0L=(0.04365)\frac{2\pi}{\SI{633e-9}{m}}(\SI{0.01}{m})=\SI{4.33e3}{rad}\approx689.5\ \text{ondas}.
\]
La parte fraccionaria ($0.53$ ondas, $191^\circ$) cambia con variaciones de espesor de apenas
$\lambda_0/\Delta n\approx\SI{15}{\micro\metre}$ por onda, o con la temperatura: en la práctica, la
polarización de salida de una lámina gruesa es impredecible.
\resultado{$n(45^\circ)=2.2436$; walk-off $\rho=2.30^\circ$; desplazamiento $\SI{0.40}{mm}$; retardo $\SI{4.33e3}{rad}$ ($\approx689.5$ ondas).}

\codigo{t3_p23_uniaxiales.py}
\begin{figure}[H]\centering
\includegraphics[width=0.55\linewidth]{t3_p23_superficie_k.png}
\caption{Problema 23(c). Superficie $\kv$ de un uniaxial positivo con el eje a $45^\circ$ de la normal ($n_e-n_o$
exagerado 5 veces para que se vea). $\kv$ es normal a la cara; el rayo extraordinario $\Sv$ es normal a la elipse
y se inclina hacia el eje óptico.}
\end{figure}

\idea{Snell vale para $\kv$, no para los rayos. El rayo e es normal a la superficie $\kv$:
$\tan\theta_S=(n_o^2/n_e^2)\tan\theta$. En un uniaxial positivo se acerca al eje; en uno negativo, se aleja.
En la doble refracción, el ángulo de $\kv$ con el eje depende del propio ángulo de refracción: hay que iterar
(6.3-25). Retardo máximo con $\kv\perp$ eje; un ``cuarto de onda'' admite $\Gamma=\pm\pi/2$ (mod $2\pi$).}

%---------------------------------------------------------------------
\sol{24}{Elipsoide de índices: puente hacia el ajuste de fase}{Saleh, §6.3C, ec.~6.3-15, p.~221; Fig.~6.3-12, p.~225; Boyd, §2.3, ecs.~2.3.5 y 2.3.8--2.3.11, pp.~80--84.}
\label{sol:24}

\inc{a} Para SHG eficiente hace falta $\Delta k=k_{2\omega}-2k_\omega=\frac{2\omega}{c_0}[n(2\omega)-n(\omega)]=0$, es decir,
$n(2\omega)=n(\omega)$ (Boyd, ec.~2.3.5, p.~81). En un medio isótropo hay un solo índice para cada
frecuencia. Con dispersión normal ($dn/d\omega>0$ entre $\omega$ y $2\omega$, lejos de resonancias),
$n(\omega)$ es estrictamente creciente, así que $n(2\omega)>n(\omega)$ siempre. Solo con una
resonancia entre $\omega$ y $2\omega$ (dispersión anómala) se podría compensar, al precio de mucha
absorción. La birrefringencia da la salida: la onda $2\omega$ viaja con la polarización ``lenta'' de
índice menor (aquí, extraordinaria en un negativo) y la fundamental con la ordinaria.

\inc{b} Tipo I \emph{ooe}: dos fotones ordinarios a $\omega$ dan uno extraordinario a $2\omega$. La condición es
$n_e(532,\theta)=n_o(1064)$ (Boyd, ec.~2.3.9). Con (6.3-15) (equivalente a Boyd 2.3.8),
$\frac{1}{n_o(1064)^2}=\frac{\cos^2\theta}{n_o(532)^2}+\frac{\sin^2\theta}{n_e(532)^2}$. Usando $\cos^2\theta=1-\sin^2\theta$ y
despejando:
\[
\sin^2\theta=\frac{n_o(1064)^{-2}-n_o(532)^{-2}}{n_e(532)^{-2}-n_o(532)^{-2}}
=\frac{0.365050-0.356467}{0.413295-0.356467}=\frac{0.008583}{0.056828}=0.15099,
\]
de donde $\sin\theta=0.38858$ y $\theta=22.87^\circ$ (Boyd, ec.~2.3.11). Existe solución porque
$n_e(532)=1.5555<n_o(1064)=1.6551<n_o(532)=1.6749$: la curva $n_e(532,\theta)$ baja de $1.6749$
a $1.5555$ y cruza el valor objetivo (Fig.~\ref{fig:t3_p24}).
\resultado{$\theta_{pm}=22.9^\circ$.}

\inc{c} La onda $2\omega$ es extraordinaria. Su rayo forma con el eje el ángulo $\theta_S$ con
\[
\tan\theta_S=\frac{n_o^2}{n_e^2}\tan\theta=\Big(\frac{1.6749}{1.5555}\Big)^2(0.42173)=1.15942\cdot0.42173=0.48896,
\]
es decir, $\theta_S=26.06^\circ$. El walk-off vale
\[
\rho=\theta_S-\theta=26.06^\circ-22.87^\circ=3.19^\circ=\SI{55.7}{mrad};
\]
en un cristal negativo el rayo e se \emph{aleja} del eje. Con la fórmula
$\tan\rho=\frac{n^2(\theta)}{2}\big(\frac{1}{n_e^2}-\frac{1}{n_o^2}\big)\sin2\theta$ sale lo mismo. Tras $L=\SI{5}{mm}$, el haz
armónico se separa del fundamental (ordinario, que va por $\kv$)
\[
\Delta x=L\tan\rho=(\SI{5}{mm})(0.0557)=\SI{279}{\micro\metre}.
\]
Con haces enfocados de algunos cientos de micras o menos, dejan de solaparse: el walk-off limita la
longitud útil del cristal (Boyd, p.~83).
\resultado{$\rho=3.19^\circ$ ($\SI{55.7}{mrad}$); separación $\SI{279}{\micro\metre}$ en $\SI{5}{mm}$.}

\inc{d} En $\theta=0$ ($\kv$ a lo largo del eje) los dos modos ven $n_o$, así que el armónico tiene $n_o(532)$:
\[
\Delta k=\frac{2\pi n_o(532)}{\lambda_{532}}-2\frac{2\pi n_o(1064)}{\lambda_{1064}}=\frac{2\pi}{\SI{532}{nm}}[n_o(532)-n_o(1064)]
=\frac{2\pi(0.0198)}{\SI{532e-9}{m}}=\SI{2.34e5}{m^{-1}},
\]
\[
L_c=\frac{\pi}{|\Delta k|}=\frac{\lambda_{1064}}{4[n_o(532)-n_o(1064)]}=\SI{13.4}{\micro\metre}.
\]
Más allá de $L_c$ la potencia armónica deja de crecer y vuelve a bajar (oscila con periodo $2L_c$, Boyd
ec.~2.3.2): sin ajuste de fase, un cristal de \SI{5}{mm} no rinde más que una capa de \SI{13}{\micro\metre}.
\resultado{$\Delta k=\SI{2.34e5}{m^{-1}}$, $L_c=\SI{13.4}{\micro\metre}$.}

\codigo{t3_p24_bbo.py}
\begin{figure}[H]\centering
\includegraphics[width=0.75\linewidth]{t3_p24_bbo.png}
\caption{Problema 24. Índices ordinario y extraordinario del BBO a 1064 y 532 nm contra el ángulo con el eje óptico.
El ajuste de fase tipo I (ooe) ocurre donde $n_e(532,\theta)$ (verde discontinua) corta a $n_o(1064)$ (azul continua):
$\theta_{pm}=22.9^\circ$.}
\label{fig:t3_p24}
\end{figure}

\idea{Fórmula de examen: $\sin^2\theta_{pm}=(n_o(\omega)^{-2}-n_o(2\omega)^{-2})/(n_e(2\omega)^{-2}-n_o(2\omega)^{-2})$
(negativo, tipo I). No confunda $n_e$, el valor principal, con $n_e(\theta)$, el de (6.3-15). El walk-off
es el precio del ajuste por ángulo, y es nulo solo en $\theta=90^\circ$ (ajuste no crítico, por temperatura).
$L_c=\lambda_\omega/[4\Delta n]$ es del orden de micras: sin ajuste de fase no hay SHG útil.}

%---------------------------------------------------------------------
\sol{25}{Cristal biaxial}{Saleh, §6.3A, ec.~6.3-8, p.~218; §6.3C, construcción del elipsoide, pp.~220--221; §6.3D, Fig.~6.3-9, p.~223.}
\label{sol:25}

\inc{a} En los ejes principales (Saleh, ec.~6.3-8):
\[
\frac{x^2}{n_x^2}+\frac{y^2}{n_y^2}+\frac{z^2}{n_z^2}=1,\qquad \frac{x^2}{1.7377^2}+\frac{y^2}{1.7453^2}+\frac{z^2}{1.8297^2}=1,
\]
con semiejes $n_x<n_y<n_z$: biaxial positivo.

\inc{b} Sea $\hat u=(\sin\theta,0,\cos\theta)$. El plano normal a $\hat u$ por el origen contiene al eje $y$ y a
la dirección $\hat v=(\cos\theta,0,-\sin\theta)$. Como el elipsoide es simétrico ante $y\to-y$, la elipse
de intersección tiene sus ejes a lo largo de $\hat y$ y de $\hat v$ (Saleh, pasos 1--3, p.~220):
\begin{itemize}
\item a lo largo de $\hat y$, el punto $(0,r,0)$ está en el elipsoide si $r=n_y$: modo con $\D\parallel\hat y$
e índice $n_a=n_y$, independiente de $\theta$ (es un ``ordinario'' dentro de este plano);
\item a lo largo de $\hat v$, el punto $r\hat v$ cumple $r^2\big(\frac{\cos^2\theta}{n_x^2}+\frac{\sin^2\theta}{n_z^2}\big)=1$:
\[
\frac{1}{n_b^2(\theta)}=\frac{\cos^2\theta}{n_x^2}+\frac{\sin^2\theta}{n_z^2},
\]
con $\D$ en el plano $x$--$z$, de $n_b(0)=n_x$ (cuando $\kv\parallel z$, $\D\parallel x$) a $n_b(90^\circ)=n_z$.
\end{itemize}
\resultado{$n_a=n_y=1.7453$ ($\D\parallel\hat y$) y $n_b(\theta)=\big[\cos^2\theta/n_x^2+\sin^2\theta/n_z^2\big]^{-1/2}$ ($\D$ en el plano $x$--$z$).}

\inc{c} Un eje óptico es una dirección en la que la elipse de índices es un círculo, es decir, $n_a=n_b$. En el plano
$x$--$z$, $n_b(\theta)$ crece de forma monótona de $n_x$ a $n_z$, y como $n_x<n_y<n_z$ cruza a $n_y$ exactamente
una vez en $(0,90^\circ)$, en $\theta=V_z$; por simetría, también en $-V_z$. La condición es
\[
\frac{1}{n_y^2}=\frac{\cos^2V_z}{n_x^2}+\frac{\sin^2V_z}{n_z^2}=\frac{1}{n_x^2}-\sin^2V_z\Big(\frac{1}{n_x^2}-\frac{1}{n_z^2}\Big).
\]
Despejando:
\[
\sin^2V_z=\frac{n_x^{-2}-n_y^{-2}}{n_x^{-2}-n_z^{-2}}=\frac{(n_y^2-n_x^2)/(n_x^2n_y^2)}{(n_z^2-n_x^2)/(n_x^2n_z^2)}
=\frac{n_z^2}{n_y^2}\,\frac{n_y^2-n_x^2}{n_z^2-n_x^2}
\;\Rightarrow\;\sin V_z=\frac{n_z}{n_y}\sqrt{\frac{n_y^2-n_x^2}{n_z^2-n_x^2}} .\ \blacksquare
\]
Numéricamente: $n_x^2=3.01960$, $n_y^2=3.04607$, $n_z^2=3.34780$, y
$\sin V_z=\frac{1.8297}{1.7453}\sqrt{\frac{0.02647}{0.32820}}=1.04836\cdot0.28399=0.29773$, de donde $V_z=17.32^\circ$.
Los ejes ópticos están en el plano $x$--$z$, el de los índices menor y mayor. En el plano $x$--$y$ los índices son
$n_z$ y uno entre $n_x$ y $n_y$; en el $y$--$z$, $n_x$ y uno entre $n_y$ y $n_z$. En ninguno de los dos
pueden coincidir. Los dos ejes forman $2V_z=34.6^\circ$ (Saleh, Fig.~6.3-9: las hojas de la superficie
$\kv$ se tocan en cuatro puntos, dos ejes).
\resultado{$V_z=17.3^\circ$ (los dos ejes ópticos están a $\pm17.3^\circ$ de $z$ en el plano $x$--$z$).}

\codigo{t3_p25_ktp.py}
\begin{figure}[H]\centering
\includegraphics[width=0.72\linewidth]{t3_p25_ktp.png}
\caption{Problema 25. Índices de los dos modos normales del KTP a 1064 nm para $\kv$ en el plano $x$--$z$: se cruzan en el
eje óptico, $\theta=V_z=17.3^\circ$.}
\end{figure}

\idea{Para propagar en un plano principal, uno de los modos tiene $\D$ perpendicular al plano y un índice constante,
y el otro se comporta como una onda ``extraordinaria'' con los dos índices del plano. Los ejes ópticos
están siempre en el plano de $n_{\min}$ y $n_{\max}$, donde el índice variable puede igualar a $n_{\rm int}$.}

%---------------------------------------------------------------------
\sol{26}{Actividad óptica y aislador de Faraday}{Saleh, §6.4A--B, ecs.~6.4-1 y 6.4-9, pp.~228--231; §6.6D, Fig.~6.6-5, p.~238.}
\label{sol:26}

\inc{a} La potencia rotatoria del cuarzo a 500 nm es $\rho\approx31^\circ/$mm (Saleh, p.~230), y el ángulo
rotado es $\rho d$:
\[
d=\frac{90^\circ}{31^\circ/\text{mm}}=\SI{2.90}{mm}.
\]
\emph{Ida y vuelta.} En la actividad óptica el vector de giro es $\mathbf G=\xi\kv$ (ec.~6.4-5): se invierte con
$\kv$. Respecto a su propia dirección de propagación, la onda que vuelve gira en el mismo sentido (por ejemplo,
horario visto de frente); visto desde el laboratorio, eso deshace la rotación de la ida. En ejes fijos del
laboratorio: ida $\T_{\rm rot}(90^\circ)$, vuelta $\T_{\rm rot}(90^\circ)^T=\T_{\rm rot}(-90^\circ)$, y el
total es $\T_{\rm rot}(-90^\circ)\T_{\rm rot}(90^\circ)=\mathbf I$. La onda vuelve a su estado inicial
(Saleh, Fig.~6.4-1b, p.~228): el medio es recíproco.
\resultado{$d=\SI{2.90}{mm}$; rotación neta en ida y vuelta: cero.}

\inc{b} \emph{Longitud.} $\rho=\mathfrak VB$ (ec.~6.4-9). Para $45^\circ=\pi/4$ rad:
\[
L=\frac{\pi/4}{|\mathfrak V|B}=\frac{0.7854}{(\SI{40}{rad/(T.m)})(\SI{1}{T})}=\SI{1.96e-2}{m}=\SI{19.6}{mm}.
\]
(El signo negativo de $\mathfrak V$ solo dice que el giro es levógiro respecto a $\mathbf B$; se orienta el imán
para que, en la ida, el giro vaya del eje del polarizador A al del B.)

\emph{Matrices de Jones} en ejes fijos del laboratorio. Polarizador A en $x$, rotador de Faraday de $+45^\circ$ y
polarizador B a $45^\circ$:
\[
\T_A=\begin{bmatrix}1&0\\0&0\end{bmatrix},\quad
\T_F=\T_{\rm rot}(45^\circ)=\frac{1}{\sqrt2}\begin{bmatrix}1&-1\\1&1\end{bmatrix},\quad
\T_B=\frac12\begin{bmatrix}1&1\\1&1\end{bmatrix}.
\]
\emph{Ida:} $\T_{\rm ida}=\T_B\T_F\T_A$. Primero $\T_F\T_A=\frac{1}{\sqrt2}\left[\begin{smallmatrix}1&0\\1&0\end{smallmatrix}\right]$, y luego
$\T_{\rm ida}=\frac{1}{\sqrt2}\left[\begin{smallmatrix}1&0\\1&0\end{smallmatrix}\right]$. La luz LP en $x$ sale como
$\frac{1}{\sqrt2}(1,1)^T$, con potencia $1$: pasa entera.

\emph{Vuelta:} el sentido del giro de Faraday lo fija $\mathbf B$, no $\kv$ ($\mathbf G=\gamma\mathbf B$, ec.~6.4-11).
En ejes del laboratorio la matriz es la \emph{misma}, $\T_F$, y no su transpuesta:
\[
\T_{\rm vuelta}=\T_A\T_F\T_B,\qquad \T_F\T_B=\frac{1}{\sqrt2}\cdot\frac12\begin{bmatrix}0&0\\2&2\end{bmatrix}=\frac{1}{\sqrt2}\begin{bmatrix}0&0\\1&1\end{bmatrix},\qquad
\T_A\T_F\T_B=\begin{bmatrix}0&0\\0&0\end{bmatrix}.
\]
$\T_{\rm vuelta}$ es la matriz nula: \emph{cualquier} onda que regrese queda bloqueada. B deja pasar la componente a
$45^\circ$, el rotador la lleva a $90^\circ$ (rotación acumulada) y A la absorbe. El script lo comprueba con
entradas $x$, $y$, $45^\circ$ y circular: transmitancia de vuelta 0 en todos los casos.
\resultado{$L=\SI{19.6}{mm}$. $\T_{\rm ida}=\frac{1}{\sqrt2}\left[\begin{smallmatrix}1&0\\1&0\end{smallmatrix}\right]$ (pasa toda la LP$_x$);
$\T_{\rm vuelta}=\T_A\T_F\T_B=\mathbf 0$.}

\inc{c} Con un rotador \emph{recíproco} (actividad óptica o nemático torcido), la matriz de vuelta en ejes del
laboratorio es $\T_F^T=\T_{\rm rot}(-45^\circ)$:
$\T_A\T_F^T\T_B=\frac{1}{\sqrt2}\left[\begin{smallmatrix}1&1\\0&0\end{smallmatrix}\right]$, y la onda de vuelta a
$45^\circ$ sale como $(1,0)^T$ y pasa A entera. La lámina $\lambda/2$ a $22.5^\circ$ hace lo mismo en la ida
($x\to45^\circ$), pero su matriz es simétrica, así que en la vuelta lleva $45^\circ\to0^\circ$ y A también la
deja pasar. En los dos casos la onda ``recorre al revés'' los cambios de la ida y vuelve a la fuente. Para aislar
hace falta romper la reciprocidad, y en óptica lineal pasiva eso solo lo hace el campo magnético
(efecto Faraday) (Saleh, p.~238).
\resultado{Rotador activo y lámina $\lambda/2$ son recíprocos: la vuelta deshace la ida y la luz regresa a la fuente
(transmitancia de vuelta 1). Solo el rotador de Faraday, no recíproco, aísla.}
\codigo{t3_p26_aislador.py}

\idea{En ejes fijos del laboratorio: recíproco $\Rightarrow$ $\T_{\rm vuelta}=\T^T$ (rotación $-\theta$);
Faraday $\Rightarrow$ $\T_{\rm vuelta}=\T$ (rotación $+\theta$ otra vez, total $2\theta$). Actividad óptica:
$\mathbf G\propto\kv$; Faraday: $\mathbf G\propto\mathbf B$. Un aislador de $45^\circ$ necesita
$L=\pi/(4|\mathfrak V|B)$.}

%---------------------------------------------------------------------
\sol{27}{Nemático torcido por matrices de Jones}{Saleh, §6.5, ecs.~6.5-1 a 6.5-10, pp.~232--234; Gooch y Tarry, \emph{J.\ Phys.\ D} \textbf{8}, 1575 (1975).}
\label{sol:27}

\textbf{[N]} \emph{Modelo.} El nemático torcido es un uniaxial local cuyo eje óptico gira como
$\theta=\alpha z$ (6.5-1); su retardo por unidad de longitud es $\beta=(n_e-n_o)k_0$ (6.5-2). Se divide el
espesor $d$ en $N=200$ capas de $\Delta z=d/N$. La capa $m$ es un retardador con el eje lento (eje óptico) a
$\theta_m=m\alpha\Delta z$ (6.5-3), con matriz propia (6.5-5) sin la fase global:
\[
\T_m=\R(-\theta_m)\begin{bmatrix}e^{-j\beta\Delta z/2}&0\\0&e^{j\beta\Delta z/2}\end{bmatrix}\R(\theta_m),\qquad
\T=\T_N\cdots\T_2\T_1\quad(6.5\text{-}6).
\]
Con torsión total $\alpha d=\pi/2$, la celda queda descrita por un solo parámetro adimensional:
$\beta d=(\beta/\alpha)(\alpha d)=(\pi/2)(\beta/\alpha)$, de modo que $u\equiv\beta d/\pi=(\beta/\alpha)/2$. Se
lanza LP en $x$, $\J_{\rm sal}=\T(1,0)^T$, y se calcula la fracción de potencia en $y$,
$|J_{{\rm sal},y}|^2/|\J_{\rm sal}|^2$, que es la transmitancia con un polarizador de salida en $y$.

\codigo{t3_p27_nematico.py}

\emph{Resultados.} La diferencia máxima con Gooch--Tarry en todo el rango es $2\times10^{-5}$: con 200 capas la
discretización es excelente.
\begin{center}
\begin{tabular}{cccc}
\toprule
$u=\beta d/\pi$ & $\beta/\alpha$ & Jones ($N=200$) & Gooch--Tarry\\
\midrule
0.05 & 0.1 & 0.0025 & 0.0025\\
0.5 & 1 & 0.2272 & 0.2272\\
1 & 2 & 0.6834 & 0.6834\\
$\sqrt3=1.732$ & 3.46 & 1.0000 & 1.0000\\
$\sqrt{15}=3.873$ & 7.75 & 1.0000 & 1.0000\\
10 & 20 & 0.9999 & 0.9999\\
25 & 50 & 0.9984 & 0.9984\\
\bottomrule
\end{tabular}
\end{center}
\begin{figure}[H]\centering
\includegraphics[width=0.78\linewidth]{t3_p27_nematico.png}
\caption{Problema 27. Fracción de potencia que sale LP en $y$ de un nemático torcido $90^\circ$ (entrada LP en $x$), contra
$\beta/\alpha$. Curva continua: producto de 200 matrices de Jones (6.5-6); discontinua: Gooch--Tarry. Las líneas
punteadas marcan los máximos exactos $u=\sqrt{4k^2-1}$.}
\end{figure}

\emph{Interpretación.}
\begin{itemize}
\item $\beta\ll\alpha$: la birrefringencia es tan débil que la luz no ``ve'' la torsión. Sale casi LP en $x$ y la
fracción en $y$ tiende a 0 ($\approx u^2$ para $u\ll1$; con $u=0.05$, $0.0025$).
\item $\beta\gg\alpha$ (régimen de Mauguin o adiabático): la onda sigue la torsión y sale LP en $y$. Es el límite
(6.5-10) de Saleh, $\T\approx\R(-\alpha d)\,\mathrm{diag}(e^{-j\beta d/2},e^{j\beta d/2})$, que para LP en $x$
(un modo del retardador) solo agrega una fase antes de girar $90^\circ$. La fracción tiende a 1 con una envolvente
$1-1/(1+u^2)$.
\item Gooch--Tarry, $1-\sin^2\!\big(\frac\pi2\sqrt{1+u^2}\big)/(1+u^2)$, vale exactamente 1 cuando
$\sqrt{1+u^2}=2k$, es decir, $u=\sqrt{3},\sqrt{15},\sqrt{35},\dots$ Son los ``mínimos'' de Gooch--Tarry: la fuga
con polarizadores paralelos, $\sin^2(\cdot)/(1+u^2)$, se anula. Las pantallas TN trabajan en el primero,
$u=\sqrt3$ ($\Delta n\,d=\sqrt3\lambda_0/2$), porque da la celda más delgada (respuesta más rápida) con
conversión completa. Entre esos puntos la salida es ligeramente elíptica: por eso la ``fracción en $y$'' baja a
$\approx0.88$ en $u\approx2.7$.
\end{itemize}
\resultado{La fracción LP en $y$ tiende a 1 para $\beta\gg\alpha$ (Mauguin) y a 0 para $\beta\ll\alpha$. Vale $0.23$ en
$u=0.5$, $0.68$ en $u=1$ y exactamente 1 en $u=\sqrt3$ (primer mínimo de Gooch--Tarry), $\sqrt{15}$, \ldots{}
El modelo de 200 capas reproduce Gooch--Tarry con un error $<3\times10^{-5}$.}

\idea{El nemático torcido es un rotador \emph{adiabático}: la condición es $\beta\gg\alpha$ (muchas vueltas de
retardo por cada radián de torsión), no $\beta d\gg1$ a secas. En la cascada (6.5-6) la primera capa va a la
derecha. El parámetro natural es $u=\Delta n\,d\cdot2/\lambda_0$; el primer mínimo de Gooch--Tarry,
$\Delta n\,d=\sqrt3\lambda_0/2\approx0.87\lambda_0$, es el punto de trabajo de una TN de $90^\circ$.}
